\chapter{Objetivos do Curso}
\label{chp:objetivos}

O \ppccurso destina-se a preparar engenheiros(as) para atuar de forma integrada nas áreas de instrumentação, controle, automação industrial, redes industriais e, como diferencial deste curso, robótica e inteligência artificial aplicada à automação. O profissional formado pela FEELT/UFU será capaz de agir com base em conhecimentos técnico-científicos sólidos e manter postura crítica, ética e alinhada às demandas do setor produtivo. Espera-se que demonstre iniciativa, competência técnica, responsabilidade profissional e compromisso com a segurança e a sustentabilidade dos processos industriais.

\section{Objetivo Geral}

Formar \ppctitulacao(s) aptos(as) a conceber, modelar, analisar, projetar, dimensionar, implementar, integrar e validar sistemas de controle e automação — incorporando, na ênfase deste curso (\ppcenfasecurricular), robótica e sistemas inteligentes —, considerando a complexidade técnica dos problemas enfrentados e os impactos econômicos, sociais, ambientais e éticos das soluções propostas. Essa formação amplia, por meio de fundamentação científica, matemática e metodológica própria da Engenharia, a base tecnológica compartilhada com o Curso Superior de Tecnologia em Automação Industrial na Área Básica de Ingresso (\autoref{chp:abi}; \autoref{sec:distincao_perfis}), capacitando o(a) egresso(a) a atuar de forma crítica na identificação e resolução de problemas de maior abrangência e complexidade do setor produtivo, em consonância com as transformações da Indústria 4.0 e as demandas do mercado de trabalho brasileiro.

\section{Objetivos Específicos}

Os objetivos deste curso são estabelecidos em consonância com as Diretrizes Curriculares Nacionais do Curso de Graduação em Engenharia~\cite{cne_ces_2_2019,cne_ces_1_2021}; a referência anterior às Diretrizes da Educação Profissional e Tecnológica e ao CNCST, herdada do Curso Superior de Tecnologia, foi removida por não se aplicar a um Bacharelado em Engenharia.

O \ppccurso busca:

\begin{itemize}
    \item Assegurar formação técnico-científica sólida nas tecnologias centrais da automação industrial — eletricidade, eletrônica, instrumentação, controle, CLP/SCADA e redes industriais — diretamente orientada ao exercício profissional;
    \item Promover a integração curricular entre os núcleos de formação básica, tecnológica/profissional, humanística e de integração/extensão, assegurando progressão adequada e interdisciplinaridade;
    \item Capacitar os(as) futuros(as) \ppctitulacao(s) a incorporar princípios de sustentabilidade, eficiência energética e segurança do trabalho em sua prática profissional;
    \item Garantir atualização permanente do currículo, flexibilizando os conteúdos mais sujeitos à obsolescência tecnológica por meio dos componentes curriculares optativos organizados em áreas de formação (\autoref{chp:modulos_optativos});
    \item Assegurar, por meio da organização em Área Básica de Ingresso (ABI) com o Curso Superior de Tecnologia em Automação Industrial (\autoref{chp:abi}), uma formação progressiva e flexível, na qual o estudante curse uma etapa inicial comum às duas graduações e, com base em escolha informada ao final do quinto período, opte pela trajetória que pretende concluir primeiro, com autonomia para buscar, posteriormente, a continuidade formativa e a eventual obtenção sucessiva das duas graduações;
    \item Proporcionar, nos períodos específicos do curso, percursos de aprofundamento técnico coerentes por meio de Ênfases Formativas — Máquinas Inteligentes, Acionamentos e Proteção Industrial (MIAPI); Robótica Autônoma e Sistemas de Controle (RASC); e Sistemas Embarcados, IoT e Computação Inteligente (SEICI) —, das quais o(a) estudante deverá integralizar, no mínimo, duas (\autoref{chp:modulos_optativos});
    \item Favorecer a realização de projetos integradores e atividades curriculares de extensão, visando o diálogo entre diferentes áreas de conhecimento e a comunidade;
    \item Integrar teoria e prática, conectando os conteúdos curriculares às realidades da atuação profissional, por meio de aulas práticas, projetos, estudos de caso e experiências aplicadas em laboratório;
    \item Reconhecer e incorporar ao percurso formativo as competências desenvolvidas em contextos não formais de aprendizagem (estágio, monitorias, extensão, vivência profissional);
    \item Acompanhar as transformações do mercado de trabalho e da indústria, fomentando respostas éticas, criativas e socialmente responsáveis às demandas consolidadas e emergentes da automação industrial;
    \item Assegurar a formação de profissionais capazes de:
    \begin{itemize}
        \item Compreender e atuar sobre questões técnicas, éticas, legais e de segurança do trabalho que permeiam a prática profissional em ambientes industriais automatizados;
        \item Atender de forma direta e aplicada às necessidades do setor produtivo por meio da automação de processos e sistemas;
        \item Contribuir com visão técnica e aplicada para a modernização de plantas e processos industriais;
        \item Cooperar em equipes multidisciplinares de engenharia, operação e manutenção, atuando com senso de responsabilidade;
        \item Utilizar de forma racional os recursos disponíveis, com foco na sustentabilidade, na segurança e na eficiência dos processos automatizados.
    \end{itemize}
\end{itemize}
